\chapter{Magnetism and Spin}\label{ch:magnetism-spin}
\index{magnetism}\index{spin}\index{Stern--Gerlach}\index{diamagnetism}\index{gyromagnetic ratio}\index{Einstein--de Haas effect}\index{quantised circulation}\index{Land\'e g-factor@Land\'e $g$-factor}

\begin{quote}
\itshape A magnetic moment is charge in motion; the only question is whether the motion is real.\\
\normalfont\hfill --- working note, 2026.
\end{quote}
\bigskip

The Schr\"odinger--Maxwell coupling of Chapter~\ref{ch:sm-coupling} gave each electron a real current
$\mathbf{j}_i = \rho_i \mathbf{v}_i$ on physical $\mathbb{R}^3$, with $\mathbf{v}_i = \nabla S_i / m_i$ the
velocity read from the phase of the complex form $\psi_i = R_i e^{iS_i}$. Magnetism is what that current does in a
magnetic field, and it is the arena where the reality of the RealQM wave function is most sharply at stake ---
because a magnetic moment is, quite literally, charge in motion. This chapter asks how much of magnetism follows
from taking the electron to be a real charge density in three-dimensional space, and the answer comes in three
graded registers: orbital magnetism, which is native and needs no complex numbers; diamagnetism, which affords a
quantitative test against experiment; and spin, which reduces to a single geometric residue.

\section{Atomic magnetism in standard quantum mechanics}\label{sec:standard-magnetism}
\index{Land\'e g-factor@Land\'e $g$-factor}\index{Hund's rules}\index{spin--orbit coupling}\index{exchange interaction}\index{Zeeman effect}

To see what RealQM adds, and what it merely reproduces, it helps to have the standard account in view first. Standard
quantum mechanics builds the magnetism of an atom as an \emph{algebra of operators} acting on the many-electron
wave function $\Psi$ on $3N$-dimensional configuration space. It has three ingredients.

\emph{Orbital moment.} The angular-momentum operator $\hat{\mathbf L}$ gives a moment
$\boldsymbol{\hat\mu}_L = -\mu_B\hat{\mathbf L}/\hbar$ with gyromagnetic factor $g_L = 1$; its eigenvalues
$L_z = m\hbar$ are quantised, $m\in\mathbb Z$, by the requirement that $\Psi$ be single-valued under rotation. The
associated ``current'' is a probability current, physical for one particle and otherwise a construct on the
$3N$-dimensional space.

\emph{Spin moment.} A second angular momentum $\hat{\mathbf S}$ is \emph{postulated} as an intrinsic, two-valued
property of the electron, with $\boldsymbol{\hat\mu}_S = -g_S\mu_B\hat{\mathbf S}/\hbar$ and $g_S \approx 2$. It is
not derived from any motion in space; it was forced empirically (the Stern--Gerlach two spots, the alkali spectral
doublets, the anomalous Zeeman effect) and \emph{justified after the fact} by the Dirac equation, from which
$g_S = 2$ falls out.

\emph{The Zeeman Hamiltonian.} Coupling to an external field $\mathbf B$ gives
\begin{equation}
\hat H_Z = \frac{\mu_B}{\hbar}\big(\hat{\mathbf L} + g_S\hat{\mathbf S}\big)\cdot\mathbf B
         + \frac{e^2}{8 m_e}\sum_i(\mathbf B\times\mathbf r_i)^2,
\label{eq:zeeman-ham}
\end{equation}
the last (quadratic) term being the diamagnetic contribution.

From these, the magnetism of a real atom is assembled by well-tested rules. The orbital and spin angular momenta
combine into a total $\hat{\mathbf J} = \hat{\mathbf L} + \hat{\mathbf S}$ (Russell--Saunders coupling for light
atoms); \emph{Hund's rules} select the ground term (maximise $S$, then $L$, then fix $J$); and the observable
moment is $g_J\mu_B\sqrt{J(J{+}1)}$ with the \emph{Land\'e} $g_J$-factor interpolating between the $g_L = 1$ and
$g_S = 2$ pieces. Open shells are then \emph{paramagnetic}, following the Curie law $\chi \propto J(J{+}1)/T$;
closed shells are \emph{diamagnetic}, governed by the $\langle r^2\rangle$ term of \eqref{eq:zeeman-ham}. Atomic
\emph{fine structure} comes from a spin--orbit coupling $\xi(r)\,\hat{\mathbf L}\!\cdot\!\hat{\mathbf S}$, a
relativistic correction; and \emph{collective} magnetism --- ferromagnetic ordering --- from the
\emph{exchange} interaction $-J_{ij}\,\mathbf S_i\!\cdot\!\mathbf S_j$ (the Heisenberg model), itself a
consequence of the antisymmetry of $\Psi$ together with the Coulomb repulsion.

This apparatus is quantitatively successful, and none of it is in dispute here: Land\'e factors, Hund's rules,
Curie and Langevin susceptibilities, fine structure and the exchange model are exactly the phenomenology RealQM
must reproduce, and largely does. Two features of the account are worth marking, because they are where RealQM
departs. First, \emph{spin is posited}, not explained: it is an intrinsic property with no picture of what is
moving, and its $g = 2$ is imported from relativity. Second, the whole construction lives in \emph{operators on
a configuration-space amplitude}, not in things in ordinary space. The remainder of this chapter re-derives the
same phenomenology as the mechanics of \emph{real circulating charge} in three-dimensional space, and asks how
far spin, too, can be reduced from a primitive to a geometric residue.

\section{Orbital magnetism is real, and needs no complex plane}\label{sec:orbital-magnetism}

It is often said that magnetism forces the theory to be complex: a circulating current needs a phase, and a phase
means a complex wave function. The truth is more precise, and more favourable to RealQM. From the density
\emph{alone} --- a single real field $\psi = R$ --- no magnetism can come, because the current
\[
\mathbf{j} = \frac{q}{m}\,\operatorname{Im}(\psi^* \nabla\psi)
\]
vanishes identically for real $\psi$: a lone amplitude carries no motion. Something beyond the bare density is
needed. But complex numbers are not the only way to carry it. Continuum mechanics never uses them, yet always
carries \emph{two} real fields, a density $\rho$ and a velocity $\mathbf{v}$. Carried the same way here --- as
they already are in Chapter~\ref{ch:sm-coupling} --- the whole of orbital magnetism is real. The current
$\mathbf{j} = \rho\mathbf{v}$ is a real circulating flow; it costs a real kinetic energy
\begin{equation}
T_{\text{flow}} = \int \tfrac{1}{2}\rho\,|\mathbf{v}|^2 \, d^3x = \int \frac{|\mathbf{j}|^2}{2\rho}\, d^3x ;
\label{eq:ms-Tflow}
\end{equation}
it carries a real angular momentum $L_z = \int \rho\,(\mathbf{x}\times\mathbf{v})_z\, d^3x$ and the ordinary
magnetostatic moment $\boldsymbol{\mu} = \tfrac{1}{2}\int \mathbf{x}\times\mathbf{j}\, d^3x$; and the magnetic
state is the stationary point of the static RealQM energy augmented by \eqref{eq:ms-Tflow} at fixed $L_z$ (or, in
an applied field, of $E - \boldsymbol{\mu}\cdot\mathbf{B}$). No $i$ enters any of this. The complex
$\psi = \sqrt{\rho}\,e^{iS}$ is exactly these two real fields repackaged into one symbol, the phase $S$ being the
velocity potential; the ``$i$'' is bookkeeping, and the physics is the flow.

\subsection*{Quantised circulation and the moment}

A stationary state with a spatial phase winding $\psi = R(\mathbf{x})\,e^{im\phi}$ carries an azimuthal current
and a magnetic moment
\begin{equation}
\mu_z = -m\,\mu_B, \qquad m \in \mathbb{Z},
\label{eq:ms-mu}
\end{equation}
quantised in the winding number and independent of the density's detailed shape: the moment-arm ($\propto s$, the
distance from the axis) and the winding velocity ($\propto 1/s$) cancel, which is the gyromagnetic relation
$\boldsymbol{\mu} = (q/2m_e)\mathbf{L}$ with gyromagnetic factor $g = 1$.

The integer $m$ is the one genuinely quantum input, and it is worth seeing that it is \emph{topological}, not a
feature of the complex plane. Single-valuedness of $e^{iS}$ --- the phase is an angle, defined modulo $2\pi$ ---
forces the circulation to close on itself,
\begin{equation}
\oint \mathbf{v}\cdot d\boldsymbol{\ell} = \frac{2\pi m}{m_e}, \qquad m\in\mathbb{Z} .
\label{eq:ms-circ}
\end{equation}
A fractional winding would be a torn, discontinuous configuration --- a branch cut in the density --- which a real
continuous charge fluid cannot be. This is precisely the Onsager--Feynman quantisation of vortices in a real
superfluid: one cannot make a half-strength vortex, because the flow must match up after a loop. What is
unavoidable here is the \emph{circle} --- the periodicity of an angle --- not the complex plane as such. The
density must also vanish on the axis for $m\neq 0$ (a hollow vortex core), since $\mathbf{v}\sim 1/s$ would
otherwise diverge.

\subsection*{Numerical check: the Zeeman response}

A Cartesian $(u,v) = (\operatorname{Re}\psi, \operatorname{Im}\psi)$ leapfrog for one electron in a
two-dimensional harmonic well, initialised in the $m = 1$ state $\psi = (x+iy)e^{-r^2/2}$, conserves norm to
$1.000000$ and holds $\mu_z/\mu_B = -1.00$ (grid value $-0.997$) over thousands of steps: the circulating density
is genuine. Adding a uniform field by minimal coupling $\mathbf{p}\to\mathbf{p} - q\mathbf{A}$ gives the Zeeman
response of Table~\ref{tab:ms-zeeman}: the $m = \pm 1$ states split by $2\mu_B B$ to four digits, while the
non-circulating $m = 0$ state is inert to linear order. This is ordinary orbital magnetism, from a charge density
in real space, with no spin invoked.

\begin{table}[h]
\centering
\begin{tabular}{cccc}
\toprule
$B$ & $E(m{=}0)$ & $E(+1)-E(-1)$ & $2\mu_B B$ \\
\midrule
$0.05$ & $0.999$ & $+0.0499$ & $0.0500$ \\
$0.10$ & $1.000$ & $+0.0999$ & $0.1000$ \\
$0.20$ & $1.004$ & $+0.1998$ & $0.2000$ \\
\bottomrule
\end{tabular}
\caption{Zeeman response of the circulating density (minimal coupling, atomic units). The $m = \pm 1$ splitting
matches $2\mu_B B$; the $m = 0$ state is inert to linear order.}
\label{tab:ms-zeeman}
\end{table}

\section{Diamagnetism: a test against experiment}\label{sec:diamagnetism}

A magnetic prediction can be read straight from the ground-state density, with no field applied at all. The
Langevin--Larmor molar diamagnetic susceptibility is
\begin{equation}
\chi_M = -\frac{N_A e^2}{6 m_e c^2}\sum_i \langle r_i^2\rangle
       = -0.792\times 10^{-6}\Big(\textstyle\sum_i \langle r_i^2\rangle / a_0^2\Big)\ \mathrm{cm^3/mol},
\label{eq:ms-langevin}
\end{equation}
a sum over electrons of the mean-square radius --- precisely the real-space quantity RealQM supplies as
$\langle r^2\rangle = \int r^2\rho\, d^3x$. We evaluate it for the closed-shell atoms He, Ne, Ar from a radial
RealQM self-consistent field of the kind used in Part~\ref{part:foundations-revisited}: each electron moves in
the kernel $-Z/r$ plus the Hartree field of \emph{all the other} electrons, with no self-repulsion (the RealQM
rule) and no exchange. The mean-square sizes feed \eqref{eq:ms-langevin} directly, giving
Table~\ref{tab:ms-dia}.

\begin{table}[h]
\centering
\begin{tabular}{lcccc}
\toprule
atom & $\sum\langle r^2\rangle\ (a_0^2)$ & $\chi_{\text{calc}}$ & $\chi_{\text{exp}}$ & ratio \\
\midrule
He & $2.38$  & $-1.88$ & $-1.90$ & $0.99$ \\
Ne & $11.14$ & $-8.82$ & $-7.20$ & $1.23$ \\
Ar & $31.66$ & $-25.1$ & $-19.6$ & $1.28$ \\
\bottomrule
\end{tabular}
\caption{Molar diamagnetic susceptibility ($10^{-6}\,\mathrm{cm^3/mol}$) from the RealQM radial SCF via
\eqref{eq:ms-langevin}, against tabulated experiment.}
\label{tab:ms-dia}
\end{table}

For helium the prediction is essentially exact ($0.99\times$): the two-electron no-self-interaction field gives
$\langle r^2\rangle = 1.19\,a_0^2$ per $1s$ electron, in agreement with Hartree--Fock and with the measured
susceptibility. This is a magnetic observable \emph{computed} from a RealQM density and matched to measurement,
not asserted. For neon and argon the size and trend are right but the atoms come out some $20$--$30\%$ too
diamagnetic; the per-shell breakdown pins the cause on the valence ($2p$ for Ne, $3p$ for Ar) being too diffuse.
Two caveats frame this honestly. First, the computation is a \emph{radial mean field}: it enforces the
no-self-interaction rule but spherically averages over the non-overlap geometry and omits exchange, both of which
would \emph{contract} the valence --- so this is closer to a lower bound on the full three-dimensional solver than
to its verdict. Second, the trend --- excellent for the two-electron atom, degrading for the many-electron ones
absent the full geometry --- is the same regime boundary RealQM meets elsewhere in this book. It is a genuine
partial success: one magnetic number secured against experiment, the right order and trend for the others, and
the missing physics named rather than hidden.

\section{Spin: the single geometric residue}\label{sec:spin-residue}

Two logically distinct things travel under the name \emph{spin}. The first is fermionic \emph{exclusion}, which
RealQM realises geometrically by non-overlapping domains with no spin variable, and which accounts --- as the
validation chapters showed --- for shells, the periodic table, and bonding. The second is the two-valued
\emph{magnetic moment}, with gyromagnetic factor $g = 2$, which the scalar theory does not carry. This leaves
exactly one residue: the electron's spin moment is $g = 2$, twice the orbital value per unit angular momentum,
and its experimental face is the Stern--Gerlach experiment.

\subsection*{Stern--Gerlach: two with no middle, and the spread}

A beam of silver atoms --- each with a single valence electron in an $L = 0$, orbitally spherical, state ---
passes through an inhomogeneous field and splits into \emph{two} spots. This is the signature of a two-valued
moment $\mu_z = \pm\mu_B$, that is spin $\tfrac{1}{2}$ with $g = 2$. An orbital moment would split into an
\emph{odd} number $2l+1$ of spots, and for the $L = 0$ silver electron would not split at all: the scalar winding
ladder $m = 0, \pm 1, \dots$ always \emph{includes zero}, an undeflected middle beam. Stern--Gerlach gives two
spots \emph{with no middle}; the missing zero is the half-integer structure $m_s = \pm\tfrac{1}{2}$, which a
scalar density cannot produce.

The \emph{separation} of the spots, as opposed to their number, is ordinary magnetostatics and native to RealQM.
An atom of mass $M$ and moment $\mu_z$ in a gradient $G = \partial_z B_z$ deflects by
$z = \mu_z G L (D + \tfrac{1}{2}L)/M v^2$ across a magnet of length $L$ and drift $D$ at speed $v$, so the split
is
\begin{equation}
\Delta z = \frac{2\mu_B\, G\, L\, (D + \tfrac{1}{2}L)}{M v^2},
\label{eq:ms-spread}
\end{equation}
\emph{linear in the forcing} $G$ and $\propto 1/M v^2$. The forcing sets the pattern's \emph{scale},
continuously; it never changes the count (two) or the value ($\pm\mu_B$), and never fills the middle. Stronger
forcing only marches the same two dots farther apart, making the absent middle \emph{more} conspicuous. The
experiment thus factors a continuous, classical magnifier (the deflection, native to a real moment in a real
gradient) from a discrete, quantised input ($\pm\mu_B$).

\subsection*{The mechanical content is already real}

One might fear a deeper obstacle: that spin's mechanical angular momentum, the $\hbar/2$ made visible by the
Einstein--de Haas effect --- a suspended bar physically rotating as it is magnetised and its moments align --- is
a further ingredient the theory must supply beyond the moment itself. In RealQM it is not. A magnetic moment here
\emph{is} a real circulating charge, and a real current carries real mechanical angular momentum by construction:
the same $L_z$ that accompanies the orbital moment in \eqref{eq:ms-mu}. So the reciprocal transfer of angular
momentum between the aligned moments and the body --- Einstein--de Haas, and its converse the Barnett effect ---
is automatic, ordinary conservation for a charge in motion, requiring no spin primitive. Everything
\emph{mechanical} about spin is already real in RealQM.

\subsection*{The quantised magnitude is topological closure}

Why is the moment quantised to $\pm\mu_B$ rather than taking an arbitrary magnitude? For the same reason the
orbital moment is. A real, continuous, single-valued configuration can only wind an integer number of times, by
the closure condition \eqref{eq:ms-circ}; a fractional moment would be a torn configuration, which does not
exist. Any picture in which a field ``steers'' the moment into two orientations is therefore constrained: the two
attractors are not free but \emph{are} the topologically allowed windings, and the continuum between them is not
merely unstable but disallowed. For spin the same closure applies to a \emph{spinorial} orientation rather than a
scalar phase: a spinor closes only after $4\pi$ (the SU(2) double cover of the rotation group), so its winding is
quantised in \emph{half}-integers, giving the two-valued $\mu_z = \pm\mu_B$ with $g = 2$. Quantised magnitude is
thus not an added postulate but the price of the reality and continuity of the charge.

\subsection*{$g = 2$ is geometric, not relativistic}

It is tempting to conclude that $g = 2$ forces relativity into the theory, since in the textbook account it falls
out of the Dirac equation. By Lévy-Leblond's theorem (1967) it does not: $g = 2$ follows already from a
first-order, \emph{non-relativistic}, Galilean spinor equation. Writing the kinetic term as
$(\boldsymbol{\sigma}\cdot\mathbf{p})^2/2m$ --- legal for free motion since
$(\boldsymbol{\sigma}\cdot\mathbf{p})^2 = p^2$ --- minimal coupling in a field gives
\begin{equation}
\frac{(\boldsymbol{\sigma}\cdot(\mathbf{p} - q\mathbf{A}))^2}{2m}
= \frac{(\mathbf{p} - q\mathbf{A})^2}{2m} - \frac{q}{2m}\,\boldsymbol{\sigma}\cdot\mathbf{B},
\label{eq:ms-pauli}
\end{equation}
the last term being exactly the $g = 2$ spin--Zeeman coupling, appearing automatically. The ``$2$'' is the SU(2)
double cover: a spinor rotates at half angle, which reciprocally doubles the moment --- a geometric, scale-free
fact, owing nothing to Lorentz invariance.

\subsection*{What is delivered, and what remains open}

The single-electron step has been carried out numerically. On a two-dimensional grid a two-component spinor
evolved with the $\boldsymbol{\sigma}\cdot\mathbf{D}$ structure ($\mathbf{D} = \mathbf{p} - q\mathbf{A}$)
reproduces the operator identity $(\boldsymbol{\sigma}\cdot\mathbf{D})^2 = \mathbf{D}^2\mathbb{1} -
q\,\boldsymbol{\sigma}\cdot\mathbf{B}$ to grid accuracy, so the $g = 2$ term \emph{emerges} from squaring rather
than being inserted by hand; and an $L = 0$ state --- the very one a scalar density leaves undeflected --- splits
into two levels $\pm\mu_B B$ with no middle. That is Stern--Gerlach, non-relativistically.

We state plainly what remains a target. First, the \emph{half-integer closure} --- that the electron's own domain
carries a spinorial ($4\pi$) rather than scalar ($2\pi$) orientation --- is exhibited by the
$\boldsymbol{\sigma}\cdot\mathbf{p}$ structure but not yet \emph{forced} by the domain geometry. Second,
\emph{closed-shell pairing}: a filled shell must pair to zero net moment (diamagnetism), but the spin-blind
non-overlap geometry does not supply antiparallel pairing --- a naive per-electron spinor leaves the two spins
degenerate and aligns them with any field (paramagnetic), and forcing continuity of the spinor across a shared
boundary costs a spin domain-wall energy that \emph{favours} the parallel configuration, the opposite of what is
needed. Closed-shell and collective magnetism therefore require an added exchange term; this is a genuine
limitation of the spin-blind geometry, of a piece with the boundary discussed in Chapter~\ref{ch:limits}.
Beyond a single splitting, the measurement statistics across incompatible Stern--Gerlach axes --- the
non-commutativity revealed by sequential analysers --- is the deeper frontier.

\section{What is new, what is not, and the boundary}\label{sec:magnetism-ledger}

Because much of the phenomenology is standard, it is worth stating plainly what RealQM does and does not add.

\emph{New is the ontology.} Magnetism becomes a literal charge fluid $(\rho, \mathbf{v})$ circulating in ordinary
space --- one genuine current per electron for any $N$ --- rather than an operator on a configuration-space
amplitude; the complex form is bookkeeping; only the integer (orbital) or half-integer (spin) winding is quantum.
It is located precisely against radiation, treated in the neighbouring chapters: both are current phenomena, but
magnetism is a \emph{spatial} winding (steady, real flow in disguise) whereas radiation is a \emph{temporal}
oscillation that genuinely requires the complex time-dependent form.

\emph{Not new is the phenomenology.} RealQM \emph{reproduces}, it does not rewrite, the gyromagnetic relation, the
quantised circulation, the Zeeman splitting, and Langevin diamagnetism. The hydrodynamic $(\rho, \mathbf{v})$
form is Madelung's (1927); quantised vortices are Onsager's and Feynman's. The novelty is in \emph{what these
quantities are}, not in their values.

\emph{The boundary is collective magnetism.} Exchange-driven magnetism --- the singlet--triplet splitting,
ferromagnetic ordering --- requires spin exchange that geometric non-overlap does not deliver, as the
closed-shell test above makes concrete. Single-electron spin is reached as a residue; the collective sector is
not.

In sum, magnetism in RealQM comes in three registers: orbital magnetism is real continuum mechanics, with the
complex plane a convenience and the integer winding the only quantum; diamagnetism is testable and, for helium,
tested and matched; and spin is a single geometric residue --- not a mechanical mystery, not an arbitrary
magnitude, not a relativistic quantity --- delivered for one electron and open for many. Nothing called spin need
be \emph{posited}: its exclusion role is geometry, its angular momentum is real charge motion, and what remains is
a single geometric factor to be \emph{derived}.
